\documentclass[12pt]{article}
\usepackage[utf8]{inputenc}
\usepackage[T1]{fontenc}
\usepackage{hyperref}
\title{Uya Language Resources for K-12 Curriculum Development}
\author{Compiled by Claude Sonnet 5 (claude-sonnet-5) for Andr\'as Kornai}
\date{}
\begin{document}
\maketitle

\section{Introduction}
Uya (ISO 639-3: usu; also Usu) is a Rai Coast language of Madang Province, Papua New Guinea, in the Trans-New Guinea family. About 270 speakers (SIL, 2003) are reported; a UNESCO-derived figure of 93 appears in the Endangered Languages Project listing. Caution: the name Uya is also used for Sausi (ssj), a different language; do not conflate them.

\section{Primary Educational Resources}
None found.

\section{Linguistic and Reference Materials}
\begin{itemize}
\item Wikipedia, Uya language: \url{https://en.wikipedia.org/wiki/Uya_language}
\item Wikipedia, Sausi language (a distinct language sometimes called Uya): \url{https://en.wikipedia.org/wiki/Sausi_language}
\end{itemize}
No dictionary, grammar or orthography was found.

\section{Religious and Community Resources}
\begin{itemize}
\item Joshua Project, Sausi/Uya people group profile (about 200 people; relation to usu not verified): \url{https://joshuaproject.net/people_groups/14789}
\end{itemize}

\section{Digital and Audio Resources}
\begin{itemize}
\item Endangered Languages Project entry (endangered): \url{https://endangeredlanguages.com/lang/3409}
\end{itemize}

\section{Curriculum Development Context}
With very few speakers and little documentation online, a curriculum would depend on new documentation and community involvement.
\end{document}